\section*{الفصل \ref{ch:g11:reaction-energy} --- طاقة التفاعلات: الاحتراق وطاقات الروابط}

\begin{solution}{exo:g11:reaction-energy:1}
\omterm{def:g4:what-a-fire-needs:burning}{احتراق} الخشب: \omterm{def:g11:reaction-energy:exothermic}{ناشر للحرارة}. انصهار الجليد: \omterm{def:g11:reaction-energy:exothermic}{ماص للحرارة} (فهو يأخذ الحرارة من
محيطه، وإن لم يكن تفاعلًا كيميائيًا). كيس التبريد:
\omterm{def:g11:reaction-energy:exothermic}{ماص للحرارة}. مدفأة اليد: ناشرة للحرارة. صنع السكر في ضوء الشمس:
\omterm{def:g11:reaction-energy:exothermic}{ماص للحرارة} (فالطاقة تأتي من الضوء).
\end{solution}

\begin{solution}{exo:g11:reaction-energy:2}
$Q = 3.0 \times 890.7 = \qty{2.7e3}{kJ}$، أي نحو \qty{2.7}{MJ}.
\end{solution}

\begin{solution}{exo:g11:reaction-energy:3}
\omterm{def:g7:chemical-reactions:reactant}{المتفاعلات}. ويُطلق الفرق إلى المحيط، غالبًا على هيئة
حرارة.
\end{solution}

\begin{solution}{exo:g11:reaction-energy:4}
$n = 100 / 46.0 = \qty{2.17}{mol}$؛ $Q = 2.17 \times 1367.6 =
\qty{2.97e3}{kJ}$، أي نحو \qty{3.0}{MJ}.
\end{solution}

\begin{solution}{exo:g11:reaction-energy:5}
الطاقة اللازمة لكسر \omterm{def:g10:the-mole:amount}{مول} واحد من هذه الروابط، \omterm{def:g7:atoms-and-molecules:molecule}{والجزيئات} \omterm{def:g7:atoms-and-molecules:atom}{والذرات}
غازية. فالزوج المشترك يمسك الذرتين معًا: وفصلهما
يعاكس هذا التجاذب، وهذا يتطلب طاقة.
\end{solution}

\begin{solution}{exo:g11:reaction-energy:6}
المكسورة: $2 \times 436 + 498 = \qty{1370}{kJ}$؛ المتكوّنة:
$4 \times 464 = \qty{1856}{kJ}$؛ $E_r \approx 1370 - 1856 =
\qty{-486}{kJ}$: \omterm{def:g11:reaction-energy:exothermic}{ناشر للحرارة}.
\end{solution}

\begin{solution}{exo:g11:reaction-energy:7}
للإيثانول \ce{CH3-CH2-O-H} خمس روابط \ce{C-H}، ورابطة \ce{C-C} واحدة، ورابطة \ce{C-O} واحدة، ورابطة
\ce{O-H} واحدة. المكسورة: $5 \times 415 + 345 + 350 + 464 + 3 \times 498 =
\qty{4728}{kJ}$. المتكوّنة: $4 \times 741 + 6 \times 464 = \qty{5748}{kJ}$.
$E_r \approx \qty{-1020}{kJ}$، أي أصغر بالقيمة المطلقة بنحو الربع من
القيمة المقيسة \qty{1367.6}{kJ}.
\end{solution}

\begin{solution}{exo:g11:reaction-energy:8}
البروبان: $2219.2 / 0.0440 = \qty{5.04e4}{kJ/kg} = \qty{50.4}{MJ/kg}$.
الميثان: $890.7 / 0.0160 = \qty{55.7}{MJ/kg}$. والقيمتان كما في المخطط.
\end{solution}

\begin{solution}{exo:g11:reaction-energy:9}
$Q = 1500 \times 4.18 \times 85 = \qty{5.33e5}{J} = \qty{533}{kJ}$؛
$n = 533 / 890.7 = \qty{0.598}{mol}$، أي $0.598 \times 16.0 =
\qty{9.6}{g}$ من الميثان؛ وإذا ضاع نصف الطاقة، \qty{19}{g}.
\end{solution}

\begin{solution}{exo:g11:reaction-energy:10}
الترتيب تنازليًا: الهيدروجين، فالميثان، فالبروبان، فالأوكتان، فالإيثانول.
$48.0 / 142.9 = \qty{0.34}{kg}$ من الهيدروجين.
\end{solution}

\begin{solution}{exo:g11:reaction-energy:11}
المكسورة: $946 + 3 \times 436 = \qty{2254}{kJ}$؛ المتكوّنة:
$6 \times 390 = \qty{2340}{kJ}$؛ $E_r \approx \qty{-86}{kJ}$: \omterm{def:g11:reaction-energy:exothermic}{ناشر للحرارة}
قليلًا.
\end{solution}

\begin{solution}{exo:g11:reaction-energy:12}
الماء: $250 \times 4.18 \times 20 = \qty{2.09e4}{J} = \qty{20.9}{kJ}$.
الإيثانول: $1.20 / 46.0 = \qty{0.0261}{mol}$، يستطيع إطلاق
$0.0261 \times 1367.6 = \qty{35.7}{kJ}$. المردود
$20.9 / 35.7 \approx \qty{59}{\percent}$. الضياع: حرارة يحملها
\omterm{def:g4:air-a-mixture-of-gases:air}{الهواء}، وحرارة تسخّن العلبة والحامل، \omterm{def:g8:combustion-and-fuels:complete}{واحتراق غير تام} (سُخام)،
وبعض الإيثانول يتبخر دون أن يحترق.
\end{solution}

\begin{solution}{exo:g11:reaction-energy:13}
\ce{2C8H18 + 25O2 -> 16CO2 + 18H2O}: 8 مولات من ثنائي أكسيد الكربون،
\qty{352}{g}، لكل \qty{5470}{kJ}، إذن $352 / 5.470 = \qty{64.4}{g}$ لكل
ميغاجول. \ce{C3H8 + 5O2 -> 3CO2 + 4H2O}: \qty{132}{g} لكل
\qty{2219.2}{kJ}، إذن \qty{59.5}{g} لكل ميغاجول. والميثان، بمقدار
\qty{49.4}{g}، يطلق أقل قدر.
\end{solution}

\begin{solution}{exo:g11:reaction-energy:14}
يعطي الجدول طاقات روابط متوسطة، في حين أن روابط \ce{C=O} في
ثنائي أكسيد الكربون أقوى من المتوسط؛ والقيم المقيسة تخص
الماء السائل، الذي يطلق تكاثفه طاقة أكثر من تفاعل
الغازات؛ والطريقة تعامل كل رابطة كأنها مستقلة عن
جاراتها. ويبقى التقدير مفيدًا: فهو يعطي الإشارة الصحيحة
(\omterm{def:g11:reaction-energy:exothermic}{ناشر للحرارة}) ورتبة المقدار الصحيحة دون أي قياس.
\end{solution}

\begin{solution}{exo:g11:reaction-energy:15}
المكسورة: $4 \times 464 = \qty{1856}{kJ}$؛ المتكوّنة: $2 \times 436 + 498 =
\qty{1370}{kJ}$؛ $E_r \approx \qty{+486}{kJ}$: \omterm{def:g11:reaction-energy:exothermic}{ماص للحرارة}، فيجب
إمداد الطاقة. وحرق الهيدروجين لاحقًا يعيد هذه الطاقة:
فالهيدروجين يخزّن طاقة أُنتجت في مكان آخر، ولا يخلق شيئًا منها.
\end{solution}

\begin{solution}{pb:g11:reaction-energy:1}
\textbf{1.} \ce{CH4 + 2O2 -> CO2 + 2H2O}.

\textbf{2.} \ce{C2H6O + 3O2 -> 2CO2 + 3H2O}.

\textbf{3.} \ce{2C8H18 + 25O2 -> 16CO2 + 18H2O}.

\textbf{4.} \ce{2H2 + O2 -> 2H2O}.

\textbf{5.} الهيدروجين.

\textbf{6.} \qty{16.0}{g/mol}، \qty{46.0}{g/mol}، \qty{114.0}{g/mol}،
\qty{2.0}{g/mol}.

\textbf{7.} بالكيلوجول لكل كيلوغرام، ثم بالوحدة \si{MJ/kg}: الميثان
$890.7 / 0.0160$، \qty{55.7}{MJ/kg}؛ الإيثانول $1367.6 / 0.0460$،
\qty{29.7}{MJ/kg}؛ الأوكتان $5470 / 0.1140$، \qty{48.0}{MJ/kg}؛ الهيدروجين
$285.8 / 0.0020$، \qty{142.9}{MJ/kg}.

\textbf{8.} الترتيب تنازليًا: الهيدروجين، فالميثان، فالأوكتان، فالإيثانول.

\textbf{9.} الهيدروجين غاز خفيف جدًا: فالكيلوغرام منه يملأ حجمًا
هائلًا تحت الضغط العادي، فيجب ضغطه في خزانات ثقيلة
عالية الضغط.

\textbf{10.} الميثان: 4 \ce{C-H}؛ ثنائي الأكسجين: 2 \ce{O=O} مكسورة. ثنائي
أكسيد الكربون: 2 \ce{C=O}؛ الماء: $2 \times 2 = 4$ \ce{O-H} متكوّنة.

\textbf{11.} $4 \times 415 + 2 \times 498 = \qty{2656}{kJ}$.

\textbf{12.} $2 \times 741 + 4 \times 464 = \qty{3338}{kJ}$.

\textbf{13.} $E_r \approx 2656 - 3338 = \qty{-682}{kJ}$، مقابل
\qty{-890.7}{kJ} مقيسة: أي أصغر بالقيمة المطلقة بنحو \qty{23}{\percent}.

\textbf{14.} طاقات الروابط متوسطة (والرابطة \ce{C=O} في ثنائي أكسيد الكربون
أقوى من المتوسط)؛ والماء سائل في القياس، والأنواع غازية في
التقدير.

\textbf{15.} $1000 / 890.7 = \qty{1.123}{mol}$.

\textbf{16.} $1.123 \times 44.0 = \qty{49.4}{g}$.

\textbf{17.} الإيثانول: $1000 / 1367.6 = \qty{0.731}{mol}$، تعطي
\qty{1.462}{mol} من ثنائي أكسيد الكربون، أي \qty{64.3}{g}. الأوكتان:
$1000 / 5470 = \qty{0.183}{mol}$، تعطي \qty{1.463}{mol}، أي \qty{64.4}{g}.

\textbf{18.} الميثان: فله أكبر عدد من \omterm{def:g7:atoms-and-molecules:atom}{ذرات} الهيدروجين لكل \omterm{def:g7:atoms-and-molecules:atom}{ذرة} كربون (4
مقابل 2.25 للأوكتان)، وحرق الهيدروجين يعطي طاقة دون
ثنائي أكسيد الكربون.

\textbf{19.} على طريقة صنع الهيدروجين: فإن صُنع من الماء بكهرباء
من مصادر متجددة لم يطلق تقريبًا أي ثنائي أكسيد الكربون؛ وإن صُنع من
الميثان أطلق ثنائي أكسيد الكربون في المصنع بدل
العادم.

\textbf{20.} يطلق الميثان نحو \qty{49}{g} من ثنائي أكسيد الكربون لكل
ميغاجول.
\end{solution}
